\section{The original relative de Rham product on every open}

\begin{lemma}[The original wedge rule on the whole twist]
\label{lem:native-di-all-open-wedge-leibniz}
\leanok
\lean{EnvelopeProperty.DI.CechWholeSectionLeibniz.formWedge_differential_onBasis, EnvelopeProperty.DI.CechWholeSectionLeibniz.formWedge_differential}
Let $k$ have prime characteristic $p$, let $X$ be smooth of relative
 dimension $n$ over $k$, and let $F:X\to X^{(p)}$ be its actual
relative Frobenius. Put $E^i=F_*\Omega^i_{X/k}$, with its original
coefficient action through $F^\#$. On every open $W\subset X^{(p)}$
and for every pair of whole sections $a\in E^q(W)$ and $b\in E^r(W)$,
the original global differential and exterior product satisfy
\[
 d(a\wedge b)=da\wedge b+(-1)^q a\wedge db.
\]
All degree transports retain the same original form; all terms
restrict by the original sheaf restriction maps.
\end{lemma}
\begin{proof}
On each member of the original twisted affine basis, the inverse
image has the original etale coordinate tuple. The global differential
restricts there to the intrinsic affine differential. The original
sheafification unit identifies every affine exterior section with
its original exterior Kahler section and retains both its differential
and its product. To compute the rule on those sections, expand each
as a finite sum of scalar multiples of ordered coordinate wedges.
For two such terms $a_0\omega$ and $b_0\eta$, where $\omega$ has
degree $q$, the coordinate wedges are closed. Applying
$d(a_0b_0)=b_0\,da_0+a_0\,db_0$ and moving $db_0$ past the $q$
one-forms of $\omega$ gives the displayed identity. Bilinearity
and additivity give it for every pair of sections on the basis open.
The explicit additive identification of original and pushed sections
preserves the calculation; the scalar actions continue to be the
actual ones through $F^\#$.

For any open $W$ and any point of $W$, choose a member of this same
basis contained in $W$. Naturality of the global differential and of
the entire original exterior product identifies the restrictions of
the two sides with the basis calculation just proved. Restriction
also preserves addition and the integer sign. Thus the two original
sections agree near every point. Separatedness of the original target
sheaf gives equality on $W$. For empty $W$ the same separation
argument applies with no point witnesses.
\end{proof}
